% Minimal XeLaTeX starter for final-review-slides artifacts (color-coded edition).
\documentclass[11pt,a4paper]{article}
\usepackage[margin=18mm]{geometry}
\usepackage{fontspec}
\usepackage{xeCJK}
\usepackage{amsmath,amssymb,booktabs,longtable,array,graphicx,hyperref}

% 五色记忆编码样式：定义蓝/绿/橙/红/黄标签、徽章、坑点/结论/机制框
% ============================================================
% color-style.tex — 五色记忆编码样式（复刻儿科学彩色速记讲义）
% 颜色语义：
%   蓝 = 基础概念   绿 = 必背数字   橙 = 机制逻辑
%   红 = 高频坑点  黄 = 核心结论
% 引入方式：在 \documentclass 之后、\begin{document} 之前
%           % ============================================================
% color-style.tex — 五色记忆编码样式（复刻儿科学彩色速记讲义）
% 颜色语义：
%   蓝 = 基础概念   绿 = 必背数字   橙 = 机制逻辑
%   红 = 高频坑点  黄 = 核心结论
% 引入方式：在 \documentclass 之后、\begin{document} 之前
%           % ============================================================
% color-style.tex — 五色记忆编码样式（复刻儿科学彩色速记讲义）
% 颜色语义：
%   蓝 = 基础概念   绿 = 必背数字   橙 = 机制逻辑
%   红 = 高频坑点  黄 = 核心结论
% 引入方式：在 \documentclass 之后、\begin{document} 之前
%           \input{color-style.tex}
% 依赖：xcolor、tcolorbox（breakable 库）；XeLaTeX 编译中文。
% ============================================================

\RequirePackage{xcolor}
\RequirePackage{tcolorbox}
\tcbuselibrary{breakable}

% ---- 深色（文字 / 边框）---- 与 docx / HTML 支线共用同一组 hex
\definecolor{pedsBlue}{HTML}{3A6EA5}
\definecolor{pedsGreen}{HTML}{2E9E6B}
\definecolor{pedsOrange}{HTML}{E08A3C}
\definecolor{pedsRed}{HTML}{CF5148}
\definecolor{pedsYellow}{HTML}{A8842A}

% ---- 浅色（标签底 / 框体底）----
\definecolor{pedsBlueBg}{HTML}{EAF2FA}
\definecolor{pedsGreenBg}{HTML}{E7F4EC}
\definecolor{pedsOrangeBg}{HTML}{FBEFE1}
\definecolor{pedsRedBg}{HTML}{FBECEA}
\definecolor{pedsYellowBg}{HTML}{FCF4DB}

% ---- 行内标签徽章 ----
\newcommand{\tagConcept}[1]{\colorbox{pedsBlueBg}{\textcolor{pedsBlue}{\textbf{#1}}}\hspace{3pt}}
\newcommand{\tagNumber}[1]{\colorbox{pedsGreenBg}{\textcolor{pedsGreen}{\textbf{#1}}}\hspace{3pt}}
\newcommand{\tagMechanism}[1]{\colorbox{pedsOrangeBg}{\textcolor{pedsOrange}{\textbf{#1}}}\hspace{3pt}}
\newcommand{\tagTrap}[1]{\colorbox{pedsRedBg}{\textcolor{pedsRed}{\textbf{#1}}}\hspace{3pt}}
\newcommand{\tagConclusion}[1]{\colorbox{pedsYellowBg}{\textcolor{pedsYellow}{\textbf{#1}}}\hspace{3pt}}

% ---- 必背数字徽章（行内强调数字 / 指标）----
\newcommand{\numBadge}[1]{\colorbox{pedsGreenBg}{\textcolor{pedsGreen}{\textbf{#1}}}}

% ---- 彩色小标题：\colorHead{pedsBlue}{标题文字} ----
\newcommand{\colorHead}[2]{\noindent\textcolor{#1}{\textbf{#2}}\par\smallskip}

% ---- 坑点框（红）：常见错误、易混、边界条件 ----
\newtcolorbox{trapbox}{%
  colback=pedsRedBg, colframe=pedsRed, breakable,
  boxrule=0.6pt, arc=6pt, boxsep=0pt,
  left=10pt, right=10pt, top=7pt, bottom=7pt}

% ---- 结论框（黄）：一句话主线、核心结论、记忆口诀 ----
\newtcolorbox{conclbox}{%
  colback=pedsYellowBg, colframe=pedsYellow, breakable,
  boxrule=0.6pt, arc=6pt, boxsep=0pt,
  left=10pt, right=10pt, top=7pt, bottom=7pt}

% ---- 机制链条框（橙，可选）：因果链、推导流程 ----
\newtcolorbox{mechbox}{%
  colback=pedsOrangeBg, colframe=pedsOrange, breakable,
  boxrule=0.6pt, arc=6pt, boxsep=0pt,
  left=10pt, right=10pt, top=7pt, bottom=7pt}

% ---- 文档开头的五色图例 ----
\newcommand{\colorLegend}{%
  \tagConcept{基础概念}\tagNumber{必背数字}\tagMechanism{机制逻辑}%
  \tagTrap{高频坑点}\tagConclusion{核心结论}}

% 依赖：xcolor、tcolorbox（breakable 库）；XeLaTeX 编译中文。
% ============================================================

\RequirePackage{xcolor}
\RequirePackage{tcolorbox}
\tcbuselibrary{breakable}

% ---- 深色（文字 / 边框）---- 与 docx / HTML 支线共用同一组 hex
\definecolor{pedsBlue}{HTML}{3A6EA5}
\definecolor{pedsGreen}{HTML}{2E9E6B}
\definecolor{pedsOrange}{HTML}{E08A3C}
\definecolor{pedsRed}{HTML}{CF5148}
\definecolor{pedsYellow}{HTML}{A8842A}

% ---- 浅色（标签底 / 框体底）----
\definecolor{pedsBlueBg}{HTML}{EAF2FA}
\definecolor{pedsGreenBg}{HTML}{E7F4EC}
\definecolor{pedsOrangeBg}{HTML}{FBEFE1}
\definecolor{pedsRedBg}{HTML}{FBECEA}
\definecolor{pedsYellowBg}{HTML}{FCF4DB}

% ---- 行内标签徽章 ----
\newcommand{\tagConcept}[1]{\colorbox{pedsBlueBg}{\textcolor{pedsBlue}{\textbf{#1}}}\hspace{3pt}}
\newcommand{\tagNumber}[1]{\colorbox{pedsGreenBg}{\textcolor{pedsGreen}{\textbf{#1}}}\hspace{3pt}}
\newcommand{\tagMechanism}[1]{\colorbox{pedsOrangeBg}{\textcolor{pedsOrange}{\textbf{#1}}}\hspace{3pt}}
\newcommand{\tagTrap}[1]{\colorbox{pedsRedBg}{\textcolor{pedsRed}{\textbf{#1}}}\hspace{3pt}}
\newcommand{\tagConclusion}[1]{\colorbox{pedsYellowBg}{\textcolor{pedsYellow}{\textbf{#1}}}\hspace{3pt}}

% ---- 必背数字徽章（行内强调数字 / 指标）----
\newcommand{\numBadge}[1]{\colorbox{pedsGreenBg}{\textcolor{pedsGreen}{\textbf{#1}}}}

% ---- 彩色小标题：\colorHead{pedsBlue}{标题文字} ----
\newcommand{\colorHead}[2]{\noindent\textcolor{#1}{\textbf{#2}}\par\smallskip}

% ---- 坑点框（红）：常见错误、易混、边界条件 ----
\newtcolorbox{trapbox}{%
  colback=pedsRedBg, colframe=pedsRed, breakable,
  boxrule=0.6pt, arc=6pt, boxsep=0pt,
  left=10pt, right=10pt, top=7pt, bottom=7pt}

% ---- 结论框（黄）：一句话主线、核心结论、记忆口诀 ----
\newtcolorbox{conclbox}{%
  colback=pedsYellowBg, colframe=pedsYellow, breakable,
  boxrule=0.6pt, arc=6pt, boxsep=0pt,
  left=10pt, right=10pt, top=7pt, bottom=7pt}

% ---- 机制链条框（橙，可选）：因果链、推导流程 ----
\newtcolorbox{mechbox}{%
  colback=pedsOrangeBg, colframe=pedsOrange, breakable,
  boxrule=0.6pt, arc=6pt, boxsep=0pt,
  left=10pt, right=10pt, top=7pt, bottom=7pt}

% ---- 文档开头的五色图例 ----
\newcommand{\colorLegend}{%
  \tagConcept{基础概念}\tagNumber{必背数字}\tagMechanism{机制逻辑}%
  \tagTrap{高频坑点}\tagConclusion{核心结论}}

% 依赖：xcolor、tcolorbox（breakable 库）；XeLaTeX 编译中文。
% ============================================================

\RequirePackage{xcolor}
\RequirePackage{tcolorbox}
\tcbuselibrary{breakable}

% ---- 深色（文字 / 边框）---- 与 docx / HTML 支线共用同一组 hex
\definecolor{pedsBlue}{HTML}{3A6EA5}
\definecolor{pedsGreen}{HTML}{2E9E6B}
\definecolor{pedsOrange}{HTML}{E08A3C}
\definecolor{pedsRed}{HTML}{CF5148}
\definecolor{pedsYellow}{HTML}{A8842A}

% ---- 浅色（标签底 / 框体底）----
\definecolor{pedsBlueBg}{HTML}{EAF2FA}
\definecolor{pedsGreenBg}{HTML}{E7F4EC}
\definecolor{pedsOrangeBg}{HTML}{FBEFE1}
\definecolor{pedsRedBg}{HTML}{FBECEA}
\definecolor{pedsYellowBg}{HTML}{FCF4DB}

% ---- 行内标签徽章 ----
\newcommand{\tagConcept}[1]{\colorbox{pedsBlueBg}{\textcolor{pedsBlue}{\textbf{#1}}}\hspace{3pt}}
\newcommand{\tagNumber}[1]{\colorbox{pedsGreenBg}{\textcolor{pedsGreen}{\textbf{#1}}}\hspace{3pt}}
\newcommand{\tagMechanism}[1]{\colorbox{pedsOrangeBg}{\textcolor{pedsOrange}{\textbf{#1}}}\hspace{3pt}}
\newcommand{\tagTrap}[1]{\colorbox{pedsRedBg}{\textcolor{pedsRed}{\textbf{#1}}}\hspace{3pt}}
\newcommand{\tagConclusion}[1]{\colorbox{pedsYellowBg}{\textcolor{pedsYellow}{\textbf{#1}}}\hspace{3pt}}

% ---- 必背数字徽章（行内强调数字 / 指标）----
\newcommand{\numBadge}[1]{\colorbox{pedsGreenBg}{\textcolor{pedsGreen}{\textbf{#1}}}}

% ---- 彩色小标题：\colorHead{pedsBlue}{标题文字} ----
\newcommand{\colorHead}[2]{\noindent\textcolor{#1}{\textbf{#2}}\par\smallskip}

% ---- 坑点框（红）：常见错误、易混、边界条件 ----
\newtcolorbox{trapbox}{%
  colback=pedsRedBg, colframe=pedsRed, breakable,
  boxrule=0.6pt, arc=6pt, boxsep=0pt,
  left=10pt, right=10pt, top=7pt, bottom=7pt}

% ---- 结论框（黄）：一句话主线、核心结论、记忆口诀 ----
\newtcolorbox{conclbox}{%
  colback=pedsYellowBg, colframe=pedsYellow, breakable,
  boxrule=0.6pt, arc=6pt, boxsep=0pt,
  left=10pt, right=10pt, top=7pt, bottom=7pt}

% ---- 机制链条框（橙，可选）：因果链、推导流程 ----
\newtcolorbox{mechbox}{%
  colback=pedsOrangeBg, colframe=pedsOrange, breakable,
  boxrule=0.6pt, arc=6pt, boxsep=0pt,
  left=10pt, right=10pt, top=7pt, bottom=7pt}

% ---- 文档开头的五色图例 ----
\newcommand{\colorLegend}{%
  \tagConcept{基础概念}\tagNumber{必背数字}\tagMechanism{机制逻辑}%
  \tagTrap{高频坑点}\tagConclusion{核心结论}}


\setmainfont{TeX Gyre Termes}
\setsansfont{TeX Gyre Heros}
\setCJKmainfont{Noto Serif CJK SC}
\hypersetup{colorlinks=true,linkcolor=pedsBlue,urlcolor=pedsBlue}
\setlength{\parindent}{0pt}
\setlength{\parskip}{0.45em}

\begin{document}

% ===== 标题区（替换为讲义标题与来源范围）=====
\section*{复习讲义标题}
\colorLegend

% ===== 用法示例：生成正式内容时删除本段 =====
\colorHead{pedsBlue}{板块小标题}

\tagConcept{概念}基础概念定义，保留教材原文表述与适用范围。

\tagNumber{数字}关键指标用 \numBadge{数字/单位} 标出，便于快速记忆。

\begin{mechbox}
  \tagMechanism{机制}因果链条：原因 \(\to\) 中间环节 \(\to\) 结果，按顺序写清推导。
\end{mechbox}

\begin{trapbox}
  \tagTrap{坑点}高频错误、易混概念、边界条件与反例写在这里。
\end{trapbox}

\begin{conclbox}
  \tagConclusion{结论}一句话主线、核心结论或记忆口诀写在这里。
\end{conclbox}

% 鉴别对照使用 longtable / tabular，列宽显式给出，避免窄列逐字折行。

\end{document}
